% Generated by the blueprint tool (blueprint-pdf/1) from a blueprint snapshot; edit the
% Markdown sources, not this file.  Compile with `latexmk -pdf` or by running
% `pdflatex` twice (the second run settles the contents and the page references).
% Images, if any, are read from assets/ next to this file.
\documentclass[11pt]{article}
\usepackage[T1]{fontenc}
\usepackage[utf8]{inputenc}
\usepackage{lmodern}
\usepackage{textcomp}
\usepackage[a4paper,margin=2.5cm,headheight=14pt]{geometry}
\usepackage{amsmath,amssymb,amsthm}
\usepackage{amscd}  % KaTeX's CD diagrams; part of amsmath, so always there
\usepackage{mathtools}
\usepackage{bm}
\usepackage{graphicx}
\usepackage{xcolor}
\usepackage{array,booktabs,tabularx}
\usepackage{fancyvrb}
\usepackage[normalem]{ulem}
\usepackage{microtype}
% Optional, and only with their fonts: a small TeX installation has the
% packages without the fonts, which fails at the first formula.
\IfFileExists{mathrsfs.sty}{\IfFileExists{rsfs10.tfm}{\usepackage{mathrsfs}}{}}{}
\IfFileExists{stmaryrd.sty}{\IfFileExists{stmary10.tfm}{\usepackage{stmaryrd}}{}}{}
\IfFileExists{cancel.sty}{\usepackage{cancel}}{}
\usepackage[hyphens]{url}
\usepackage[pdfusetitle,bookmarksnumbered=false]{hyperref}
\usepackage{bookmark}
\definecolor{bplink}{HTML}{1F4E9C}
\definecolor{bpmuted}{HTML}{5F676E}
\hypersetup{colorlinks=true,linkcolor=bplink,urlcolor=bplink,citecolor=bplink}
\setlength{\parindent}{0pt}
\setlength{\parskip}{0.55em plus 0.2em minus 0.1em}
% Lists as tight as the page's.
\makeatletter
\def\@listi{\leftmargin\leftmargini\parsep 0pt\topsep 0.3em plus 0.1em\itemsep 0.1em plus 0.1em}
\let\@listI\@listi
\def\@listii{\leftmargin\leftmarginii\labelwidth\leftmarginii\advance\labelwidth-\labelsep\parsep 0pt\topsep 0.1em\itemsep 0.1em}
% An ordered list that does not start at 1.
\newcommand{\bpstartat}[1]{\setcounter{\@enumctr}{\numexpr#1-1\relax}}
\makeatother
\setcounter{tocdepth}{3}
\emergencystretch=3em
\allowdisplaybreaks

% A numbered statement, set as a paper sets one: "Theorem 2.1 (Title)."
% #1 is the lead ("Theorem~2.1"), #2 what follows it in roman.
\newenvironment{bpstatement}[2]{\par\addvspace{0.6\baselineskip}\noindent\textbf{#1}#2\textbf{.}\enspace\ignorespaces}
  {\par\addvspace{0.4\baselineskip}}
\newenvironment{bpstatementit}[2]{\begin{bpstatement}{#1}{#2}\itshape\ignorespaces}{\end{bpstatement}}
% A step: a relation with prose of its own, read under the object it starts from.
\newenvironment{bpstep}[1]{\par\addvspace{0.3\baselineskip}\begin{list}{}{\leftmargin=1.6em\rightmargin=0pt\topsep=0pt}\item\relax
  {\small\color{bpmuted}#1}\par\nopagebreak}{\end{list}\par}
% Notes about where an object is written out.
\newenvironment{bpnote}{\par\small\color{bpmuted}\itshape}{\par}
\newcommand{\bpomitted}[1]{\par{\small\color{bpmuted}\itshape #1\par}}
\newcommand{\bpbodyheading}[1]{\par\addvspace{0.4\baselineskip}\noindent\textbf{#1}\par\nopagebreak}
\newcommand{\bprule}{\par\medskip\noindent\hfil\rule{0.3\linewidth}{0.4pt}\hfil\par\medskip}
\newcommand{\bpboundary}[1]{\par\nopagebreak{\small\color{bpmuted}#1\par}\nopagebreak}
% A reference to something this export leaves out, and one that resolves nowhere.
\newcommand{\bpout}[1]{#1\textsuperscript{\hyperlink{bp-outside}{\ensuremath{\ast}}}}
\newcommand{\bpbroken}[1]{\textcolor{red!70!black}{[#1?]}}
\newcommand{\bpslug}[1]{\texttt{#1}}
% A bibliography entry: label, anchor, entry; the label hangs in the margin.
\newcommand{\bpbibitem}[3]{\par\setbox0\hbox{[#1]\enspace}\ifdim\wd0<5em\setbox0\hbox to 5em{[#1]\hfil}\fi
  \hangindent=\wd0 \hangafter=1 \noindent\box0\phantomsection\label{#2}#3\par\smallskip}
% A character the export knows no way to set.
\newcommand{\bpglyph}[1]{\fbox{\scriptsize #1}}
% Images, at their own size unless they are wider than the line.
\newsavebox{\bpimagebox}
\newcommand{\bpimage}[1]{\sbox{\bpimagebox}{\includegraphics{#1}}%
  \ifdim\wd\bpimagebox>\linewidth\resizebox{\linewidth}{!}{\usebox{\bpimagebox}}\else\usebox{\bpimagebox}\fi}
\newcommand{\bpmissingimage}[1]{\fbox{\parbox{0.8\linewidth}{\centering\small\color{bpmuted}Image not included: #1}}}
\DefineVerbatimEnvironment{bpverbatim}{Verbatim}{fontsize=\small,xleftmargin=1em}
\renewcommand{\qedsymbol}{\ensuremath{\blacksquare}}
\hypersetup{pdftitle={Weil Conjectures — 18 The Deligne–Laumon theorem in general},pdfsubject={Chapter 18 The Deligne–Laumon theorem in general; all levels},pdfcreator={blueprint}}
% KaTeX commands the website knows
\providecommand{\R}{\mathbb{R}}\providecommand{\N}{\mathbb{N}}\providecommand{\Z}{\mathbb{Z}}
\providecommand{\Q}{\mathbb{Q}}\providecommand{\C}{\mathbb{C}}
\providecommand{\Reals}{\mathbb{R}}\providecommand{\reals}{\mathbb{R}}\providecommand{\Complex}{\mathbb{C}}
\providecommand{\cnums}{\mathbb{C}}\providecommand{\natnums}{\mathbb{N}}\providecommand{\Bbb}[1]{\mathbb{#1}}
\providecommand{\bold}[1]{\mathbf{#1}}\providecommand{\frak}[1]{\mathfrak{#1}}\providecommand{\mathbbm}[1]{\mathbb{#1}}
\providecommand{\lang}{\langle}\providecommand{\rang}{\rangle}\providecommand{\bra}[1]{\langle #1\rvert}
\providecommand{\ket}[1]{\lvert #1\rangle}\providecommand{\braket}[1]{\langle #1\rangle}
\providecommand{\set}[1]{\{#1\}}\providecommand{\Set}[1]{\left\{#1\right\}}
\providecommand{\infin}{\infty}\providecommand{\isin}{\in}\providecommand{\plusmn}{\pm}\providecommand{\sub}{\subset}
\providecommand{\sube}{\subseteq}\providecommand{\supe}{\supseteq}\providecommand{\exist}{\exists}
\providecommand{\bull}{\bullet}\providecommand{\larr}{\leftarrow}\providecommand{\rarr}{\rightarrow}
\providecommand{\lrarr}{\leftrightarrow}\providecommand{\Larr}{\Leftarrow}\providecommand{\Rarr}{\Rightarrow}
\providecommand{\Lrarr}{\Leftrightarrow}\providecommand{\uarr}{\uparrow}\providecommand{\darr}{\downarrow}
\providecommand{\harr}{\leftrightarrow}\providecommand{\lArr}{\Leftarrow}\providecommand{\rArr}{\Rightarrow}
\providecommand{\hArr}{\Leftrightarrow}\providecommand{\alef}{\aleph}\providecommand{\alefsym}{\aleph}
\providecommand{\real}{\Re}\providecommand{\image}{\Im}\providecommand{\weierp}{\wp}
\providecommand{\Alpha}{\mathrm{A}}\providecommand{\Beta}{\mathrm{B}}\providecommand{\Epsilon}{\mathrm{E}}
\providecommand{\Zeta}{\mathrm{Z}}\providecommand{\Eta}{\mathrm{H}}\providecommand{\Iota}{\mathrm{I}}
\providecommand{\Kappa}{\mathrm{K}}\providecommand{\Mu}{\mathrm{M}}\providecommand{\Nu}{\mathrm{N}}
\providecommand{\Omicron}{\mathrm{O}}\providecommand{\Rho}{\mathrm{P}}\providecommand{\Tau}{\mathrm{T}}
\providecommand{\Chi}{\mathrm{X}}\providecommand{\omicron}{o}\providecommand{\gt}{>}\providecommand{\lt}{<}
\providecommand{\And}{\&}\providecommand{\hdashline}{\hline}\providecommand{\operatornamewithlimits}[1]{\operatorname*{#1}}
\providecommand{\llbracket}{[\![}\providecommand{\rrbracket}{]\!]}\providecommand{\cancel}[1]{#1}
\providecommand{\mathscr}{\mathcal}
% the characters beyond ASCII this document uses
\DeclareUnicodeCharacter{2160}{\textup{I}}
\DeclareUnicodeCharacter{2161}{\textup{II}}
\DeclareUnicodeCharacter{2162}{\textup{III}}
% the project's macros ([katex.macros] in blueprint.toml)
\providecommand{\A}{}\renewcommand{\A}{\mathbf A}
\providecommand{\AEn}{}\renewcommand{\AEn}{A_{E,n}}
\providecommand{\Aaf}{}\renewcommand{\Aaf}{\mathbf A}
\providecommand{\Ab}{}\renewcommand{\Ab}{\mathsf{Ab}}
\providecommand{\AffSch}{}\renewcommand{\AffSch}{\operatorname{AffSch}}
\providecommand{\Aone}{}\renewcommand{\Aone}{\mathbf A^1}
\providecommand{\Bl}{}\renewcommand{\Bl}{\operatorname{Bl}}
\providecommand{\Br}{}\renewcommand{\Br}{\mathrm{Br}}
\providecommand{\C}{}\renewcommand{\C}{\mathbb C}
\providecommand{\Cech}{}\renewcommand{\Cech}{\check{H}}
\providecommand{\Coext}{}\renewcommand{\Coext}{\operatorname{Coext}}
\providecommand{\Cone}{}\renewcommand{\Cone}{\operatorname{Cone}}
\providecommand{\Cons}{}\renewcommand{\Cons}{\operatorname{Cons}}
\providecommand{\Cref}{}\renewcommand{\Cref}[1]{\text{#1}}
\providecommand{\D}{}\renewcommand{\D}{D}
\providecommand{\DLAFT}{}\renewcommand{\DLAFT}{\operatorname{FT}}
\providecommand{\DLAFr}{}\renewcommand{\DLAFr}{\operatorname{Fr}}
\providecommand{\DLAQ}{}\renewcommand{\DLAQ}{{\overline{\mathbf Q}_\ell}}
\providecommand{\DLRQ}{}\renewcommand{\DLRQ}{{\overline{\mathbf Q}_\ell}}
\providecommand{\Dbc}{}\renewcommand{\Dbc}{D^b_c}
\providecommand{\Dc}{}\renewcommand{\Dc}{D_{\mathrm{cons}}}
\providecommand{\Dctf}{}\renewcommand{\Dctf}{D^b_{ctf}}
\providecommand{\Dic}{}\renewcommand{\Dic}{D_{\IndCons}}
\providecommand{\Dplus}{}\renewcommand{\Dplus}{D^{+}}
\providecommand{\End}{}\renewcommand{\End}{\operatorname{End}}
\providecommand{\Et}{}\renewcommand{\Et}{\mathrm{\acute{e}t}}
\providecommand{\Ex}{}\renewcommand{\Ex}{\mathrm{Ex}}
\providecommand{\Ext}{}\renewcommand{\Ext}{\operatorname{Ext}}
\providecommand{\F}{}\renewcommand{\F}{\mathbb F}
\providecommand{\FT}{}\renewcommand{\FT}{\operatorname{FT}}
\providecommand{\Fbar}{}\renewcommand{\Fbar}{{\overline{\mathbf F}_q}}
\providecommand{\Fell}{}\renewcommand{\Fell}{\mathbf F_\ell}
\providecommand{\Fib}{}\renewcommand{\Fib}{\operatorname{Fib}}
\providecommand{\Fix}{}\renewcommand{\Fix}{\operatorname{Fix}}
\providecommand{\Fp}{}\renewcommand{\Fp}{{\mathbf F_{p}}}
\providecommand{\Fq}{}\renewcommand{\Fq}{{\mathbf F_q}}
\providecommand{\Fqn}{}\renewcommand{\Fqn}{{\mathbf F_{q^{n}}}}
\providecommand{\Frob}{}\renewcommand{\Frob}{\operatorname{Frob}}
\providecommand{\GL}{}\renewcommand{\GL}{\operatorname{GL}}
\providecommand{\Gal}{}\renewcommand{\Gal}{\operatorname{Gal}}
\providecommand{\Hens}{}\renewcommand{\Hens}{\operatorname{Hens}}
\providecommand{\Hilb}{}\renewcommand{\Hilb}{\operatorname{Hilb}}
\providecommand{\Hom}{}\renewcommand{\Hom}{\operatorname{Hom}}
\providecommand{\IndCons}{}\renewcommand{\IndCons}{\operatorname{IndCons}}
\providecommand{\Kar}{}\renewcommand{\Kar}{\operatorname{Kar}}
\providecommand{\LT}{}\renewcommand{\LT}{\operatorname{LT}}
\providecommand{\Lie}{}\renewcommand{\Lie}{\operatorname{Lie}}
\providecommand{\Loc}{}\renewcommand{\Loc}{\operatorname{Loc}}
\providecommand{\Mod}{}\renewcommand{\Mod}{\mathsf{Mod}}
\providecommand{\N}{}\renewcommand{\N}{\mathbb{N}}
\providecommand{\Needspace}{}\renewcommand{\Needspace}{\par}
\providecommand{\O}{}\renewcommand{\O}{{\mathcal{O}}}
\providecommand{\Ob}{}\renewcommand{\Ob}{\mathrm{Ob}}
\providecommand{\Oe}{}\renewcommand{\Oe}{\mathcal O_E}
\providecommand{\PP}{}\renewcommand{\PP}{\mathbf{P}}
\providecommand{\Perf}{}\renewcommand{\Perf}{\operatorname{Perf}}
\providecommand{\Pic}{}\renewcommand{\Pic}{\mathrm{Pic}}
\providecommand{\Pm}{}\renewcommand{\Pm}{\mathbf P^{m}}
\providecommand{\Pone}{}\renewcommand{\Pone}{\mathbf P^{1}}
\providecommand{\Pqn}{}\renewcommand{\Pqn}{\mathbf P^{n}}
\providecommand{\ProFin}{}\renewcommand{\ProFin}{\operatorname{ProFin}}
\providecommand{\Proet}{}\renewcommand{\Proet}{\mathrm{pro\acute{e}t}}
\providecommand{\Q}{}\renewcommand{\Q}{\mathbb Q}
\providecommand{\Qb}{}\renewcommand{\Qb}{\overline{\mathbf Q}_{\ell}}
\providecommand{\Qbar}{}\renewcommand{\Qbar}{\overline{\mathbf Q}}
\providecommand{\QbarX}{}\renewcommand{\QbarX}{\overline{\mathbf Q}_{\ell,X}}
\providecommand{\Qbarl}{}\renewcommand{\Qbarl}{{\overline{\mathbf Q}_{\ell}}}
\providecommand{\Qell}{}\renewcommand{\Qell}{{\mathbf{Q}_{\ell}}}
\providecommand{\Ql}{}\renewcommand{\Ql}{\mathbf Q_{\ell}}
\providecommand{\Qlb}{}\renewcommand{\Qlb}{{\overline{\mathbf Q}_\ell}}
\providecommand{\Qlbar}{}\renewcommand{\Qlbar}{\overline{\mathbf Q}_\ell}
\providecommand{\R}{}\renewcommand{\R}{\mathbb{R}}
\providecommand{\RG}{}\renewcommand{\RG}{\operatorname{\mathrm{R}\Gamma}}
\providecommand{\RGamma}{}\renewcommand{\RGamma}{R\Gamma}
\providecommand{\RGammac}{}\renewcommand{\RGammac}{R\Gamma_c}
\providecommand{\RHom}{}\renewcommand{\RHom}{R\mathcal{H}om}
\providecommand{\RHomX}{}\renewcommand{\RHomX}{R\mathcal H om_X}
\providecommand{\Rat}{}\renewcommand{\Rat}{\operatorname{IsRat}}
\providecommand{\Res}{}\renewcommand{\Res}{\operatorname{Res}}
\providecommand{\Rlim}{}\renewcommand{\Rlim}{R\!\varprojlim}
\providecommand{\SIm}{}\renewcommand{\SIm}{\operatorname{SIm}}
\providecommand{\Sch}{}\renewcommand{\Sch}{\mathsf{Sch}}
\providecommand{\Sh}{}\renewcommand{\Sh}{\operatorname{Sh}}
\providecommand{\Span}{}\renewcommand{\Span}{\operatorname{span}}
\providecommand{\Spec}{}\renewcommand{\Spec}{\operatorname{Spec}}
\providecommand{\Tr}{}\renewcommand{\Tr}{\operatorname{Tr}}
\providecommand{\Trsm}{}\renewcommand{\Trsm}{\operatorname{Tr}^{\mathrm{sm}}}
\providecommand{\Wt}{}\renewcommand{\Wt}{\operatorname{Wt}}
\providecommand{\Z}{}\renewcommand{\Z}{\mathbb{Z}}
\providecommand{\Zar}{}\renewcommand{\Zar}{\mathrm{zar}}
\providecommand{\Zell}{}\renewcommand{\Zell}{{\mathbf{Z}_{\ell}}}
\providecommand{\Zl}{}\renewcommand{\Zl}{{\mathbf Z_\ell}}
\providecommand{\affet}{}\renewcommand{\affet}[1]{{#1}^{\mathrm{aff}}_{\mathrm{ét}}}
\providecommand{\affproet}{}\renewcommand{\affproet}[1]{{#1}^{\mathrm{aff}}_{\mathrm{proét}}}
\providecommand{\boxt}{}\renewcommand{\boxt}{\boxtimes_A^{\mathrm L}}
\providecommand{\cF}{}\renewcommand{\cF}{\mathcal{F}}
\providecommand{\cHom}{}\renewcommand{\cHom}{R\mathcal{H}\!om}
\providecommand{\cI}{}\renewcommand{\cI}{\mathcal{I}}
\providecommand{\cL}{}\renewcommand{\cL}{\mathcal L}
\providecommand{\cM}{}\renewcommand{\cM}{\mathcal M}
\providecommand{\cN}{}\renewcommand{\cN}{\mathcal N}
\providecommand{\cO}{}\renewcommand{\cO}{\mathcal O}
\providecommand{\cT}{}\renewcommand{\cT}{\mathcal T}
\providecommand{\checked}{}\renewcommand{\checked}{\color{red}\aftergroup\restorecolor}
\providecommand{\coeff}{}\renewcommand{\coeff}{\operatorname{coeff}}
\providecommand{\coker}{}\renewcommand{\coker}{\operatorname{coker}}
\providecommand{\colim}{}\renewcommand{\colim}{\operatorname{colim}}
\providecommand{\constructedby}{}\renewcommand{\constructedby}[1]{\par\smallskip\noindent\textcolor{purple!60!black}{\textsf{Constructed by:} }
  \footnotesize\sloppy\nolinkurl{#1}\par}
\providecommand{\deepresult}{}\renewcommand{\deepresult}[1]{\par\smallskip\noindent\textcolor{orange!45!black}{
  \textsf{Blueprint role: named deep theorem node; see the cluster proof architecture.}}
  \par\noindent\textcolor{orange!45!black}{\textsf{Mathematical source:} }#1\par}
\providecommand{\dt}{}\renewcommand{\dt}{\mathbin{\otimes^{\mathbf L}}}
\providecommand{\dw}{}\renewcommand{\dw}{\operatorname{dw}}
\providecommand{\et}{}\renewcommand{\et}[1]{{#1}_{\mathrm{ét}}}
\providecommand{\ev}{}\renewcommand{\ev}{\operatorname{ev}}
\providecommand{\fib}{}\renewcommand{\fib}{\operatorname{fib}}
\providecommand{\formalizationnote}{}\renewcommand{\formalizationnote}[1]{\begin{quote}\small\raggedright\textbf{Formalization note.} #1\end{quote}}
\providecommand{\fpqc}{}\renewcommand{\fpqc}{X_{\mathrm{fpqc}}}
\providecommand{\gammacont}{}\renewcommand{\gammacont}{\Gamma_{\mathrm{cont}}}
\providecommand{\gm}{}\renewcommand{\gm}{\mathbf G_m}
\providecommand{\hocolim}{}\renewcommand{\hocolim}{\operatorname{hocolim}}
\providecommand{\holim}{}\renewcommand{\holim}{\operatorname{holim}}
\providecommand{\hrB}{}\renewcommand{\hrB}{\@ifnextchar]{\hrEnd}{\hrStep}}
\providecommand{\hrEnd}{}\renewcommand{\hrEnd}{\@firstoftwo{}}
\providecommand{\hrStep}{}\renewcommand{\hrStep}{\@firstoftwo{\hrB}}
\providecommand{\id}{}\renewcommand{\id}{\operatorname{id}}
\providecommand{\kunFl}{}\renewcommand{\kunFl}{\mathbf F_\ell}
\providecommand{\kunZn}{}\renewcommand{\kunZn}{\mathbf Z/n\mathbf Z}
\providecommand{\kuncd}{}\renewcommand{\kuncd}{\operatorname{cd}}
\providecommand{\kundim}{}\renewcommand{\kundim}{\dim}
\providecommand{\kunotimes}{}\renewcommand{\kunotimes}[1]{\otimes^{\mathbf L}_{#1}}
\providecommand{\kunsupp}{}\renewcommand{\kunsupp}{\operatorname{Supp}}
\providecommand{\lotimes}{}\renewcommand{\lotimes}{\otimes^{\mathrm L}}
\providecommand{\m}{}\renewcommand{\m}{{\mathfrak{m}}}
\providecommand{\n}{}\renewcommand{\n}{{\mathfrak{n}}}
\providecommand{\nagataGap}{}\renewcommand{\nagataGap}[1]{\par\noindent\textbf{Formalization gap.} #1\par}
\providecommand{\openobligation}{}\renewcommand{\openobligation}[1]{\par\smallskip\noindent\textcolor{red!65!black}{\textsf{Open obligation:} }
  \footnotesize\sloppy\nolinkurl{#1}\par}
\providecommand{\p}{}\renewcommand{\p}{{\mathfrak{p}}}
\providecommand{\pr}{}\renewcommand{\pr}{\operatorname{pr}}
\providecommand{\preshv}{}\renewcommand{\preshv}[1]{\mathrm{PreShv}(#1)}
\providecommand{\pro}{}\renewcommand{\pro}[1]{\mathrm{Pro}(#1)}
\providecommand{\proet}{}\renewcommand{\proet}{X_{\mathrm{proét}}}
\providecommand{\q}{}\renewcommand{\q}{{\mathfrak{q}}}
\providecommand{\rad}{}\renewcommand{\rad}{\operatorname{rad}}
\providecommand{\rank}{}\renewcommand{\rank}{\operatorname{rank}}
\providecommand{\red}{}\renewcommand{\red}{\color{red}}
\providecommand{\restorecolor}{}\renewcommand{\restorecolor}{\color{black}}
\providecommand{\retiredobligation}{}\renewcommand{\retiredobligation}[1]{\par\smallskip\noindent\textcolor{red!75!black}{\textsf{Retired invalid route:} }
  \footnotesize\sloppy\nolinkurl{#1}\par}
\providecommand{\rk}{}\renewcommand{\rk}{\operatorname{rk}}
\providecommand{\sep}{}\renewcommand{\sep}{\textnormal{sep}}
\providecommand{\shv}{}\renewcommand{\shv}[1]{\mathrm{Shv}(#1)}
\providecommand{\spec}{}\renewcommand{\spec}[1]{\mathrm{Spec}(#1)}
\providecommand{\stacks}{}\renewcommand{\stacks}[1]{\href{https://stacks.math.columbia.edu/tag/#1}{\textnormal{Stacks, Tag~#1}}}
\providecommand{\stacksproject}{}\renewcommand{\stacksproject}[1]{{\cite[\href{https://stacks.math.columbia.edu/tag/#1}{Tag #1}]{stacks-project}}}
\providecommand{\statusnote}{}\renewcommand{\statusnote}[1]{\begin{quote}\small\raggedright\textbf{Status.} #1\end{quote}}
\providecommand{\tildeideal}{}\renewcommand{\tildeideal}[2]{\widetilde{(#1, #2)}}
\providecommand{\tildering}{}\renewcommand{\tildering}[3]{{#1}_{\widetilde{(#2, #3)}}}
\providecommand{\tr}{}\renewcommand{\tr}{\operatorname{tr}}
\providecommand{\wt}{}\renewcommand{\wt}{\operatorname{wt}}

\begin{document}
\pagestyle{plain}
\begin{center}
{\LARGE\bfseries Weil Conjectures\par}
\vspace{0.7em}
{\large Chapter 18: The Deligne–Laumon theorem in general\par}
\vspace{0.5em}
{\small\color{bpmuted}Snapshot 2026-09-26, revision 8e4cb88e2a\par}
\end{center}
\vspace{0.5em}
{\small\tableofcontents}
\bigskip\hrule\bigskip

% object: chap-The-Deligne-Laumon-theorem-in-general (blueprint/chap-The-Deligne-Laumon-theorem-in-general/_section.md)
\section*{\texorpdfstring{18\quad The Deligne–Laumon theorem in general}{18 The Deligne–Laumon theorem in general}}
\phantomsection\label{bp1}%
\addcontentsline{toc}{section}{\texorpdfstring{18\quad The Deligne–Laumon theorem in general}{18 The Deligne–Laumon theorem in general}}
% text: prose-18-01
\phantomsection\label{bp2}%
At this point we know the Deligne–Laumon theorem on curves: a pure lisse sheaf of weight \(w\) on a curve has \(H^r_c\) of weights at most \(w+r\). For the Weil conjectures we need the same bound for the cohomology of a variety of any dimension. In this chapter we prove Theorem \bpout{13.11}, which says that for a separated morphism of finite type \(f:X_0\to Y_0\), the functor \(Rf_!\) maps \(D_{\leq w}(X_0)\) into \(D_{\leq w}(Y_0)\).

Take a surface fibred in curves over a curve, \(f:X_0\to Y_0\). Here the Leray spectral sequence \(H^a_c(Y,R^bf_!\overline{\mathbf Q}_\ell)\Rightarrow H^{a+b}_c(X,\overline{\mathbf Q}_\ell)\) suggests applying the curve case twice, once on the fibres and once on the base. On the fibres there is no difficulty: at a closed point \(y\) the stalk of \(R^bf_!\overline{\mathbf Q}_\ell\) is \(H^b_c(X_{\bar y},\overline{\mathbf Q}_\ell)\), and the curve case says that it has weights at most \(b\). On the base, however, the curve case applies only to pure sheaves. So we need a finite filtration of \(R^bf_!\overline{\mathbf Q}_\ell\) with pure quotients, that is, we need \(R^bf_!\overline{\mathbf Q}_\ell\) to be mixed in the sense of Definition \bpout{13.1}, and bounds at each point separately do not produce such a filtration. In his first proof [\hyperref[bp49]{Del74}] Deligne avoided the problem with a Lefschetz pencil, where the Picard–Lefschetz formula describes the direct image in the middle degree, but for an arbitrary fibration the sheaves on the base are not explicit. Once a mixed filtration exists, on the other hand, the pointwise bound also bounds the weights of its quotients (Lemma \bpout{13.3}). The question is therefore whether the direct images are mixed.

The answer comes once more from positivity (Theorem \hyperref[bp25]{18.8}):

\begin{quote}
every constructible real Weil sheaf on a separated scheme of finite type over \(k_0\) is mixed.
\end{quote}

To see why, let \(S\) be an irreducible constituent of a real lisse sheaf on a smooth affine \(X_0\). In order to compare the eigenvalues of \(S\) at a fixed point \(x_*\) with those at an arbitrary point \(x\), we restrict to a curve through both. On the curve each irreducible constituent is pure by Theorem \bpout{14.26}, so if the restriction of \(S\) stays irreducible, its eigenvalues at \(x\) have the weight read off at \(x_*\). In order to find curves on which \(S\) stays irreducible, we first show that irreducibility can be detected in one finite quotient of the fundamental group (Lemma \hyperref[bp8]{18.2}), and then Bertini's theorem gives a curve through the given points that stays connected in the corresponding finite étale cover (Lemma \hyperref[bp11]{18.3} and Corollary \hyperref[bp14]{18.4}). For constructible sheaves, finally, we stratify.

By Noether normalization over strata (Lemma \hyperref[bp42]{18.14}), an arbitrary \(f\) is built from quasi-finite maps and projections of relative affine lines. The essential case is a smooth family of affine curves \(h:U_0\to Y_0\) with a pure lisse sheaf \(F\) on \(U_0\). To apply the theorem on real sheaves we need direct images that are real, and here the curve case helps a second time. After a twist to weight zero and the addition of the dual, \(F\) becomes a real sheaf \(G\), and the curve case on the fibres makes the direct images of \(G\) real, hence mixed. A twist of \(F\) is a direct summand of \(G\), so the direct images of \(F\) are mixed as well (Theorem \hyperref[bp29]{18.9}).

% object: sec-Quasi-finite-morphisms-and-curves (blueprint/chap-The-Deligne-Laumon-theorem-in-general/sec-Quasi-finite-morphisms-and-curves/_section.md), content hidden
\subsection*{\texorpdfstring{18.1\quad Quasi-finite morphisms and curves}{18.1 Quasi-finite morphisms and curves}}
\phantomsection\label{bp3}%
\addcontentsline{toc}{subsection}{\texorpdfstring{18.1\quad Quasi-finite morphisms and curves}{18.1 Quasi-finite morphisms and curves}}
% text: prose-18-02
\phantomsection\label{bp4}%
To prove Theorem \hyperref[bp25]{18.8} we restrict a sheaf to curves, reassemble it over strata and descend it along constant extensions. For the reassembly and the descent we need that quasi-finite direct images preserve weights. For the restriction we need curves on which given lisse sheaves stay geometrically irreducible.

The first fact is elementary, because a finite direct image only permutes stalks. Let \(y\) be a point downstairs and \(x\) a point over \(y\) with residue extension of degree \(e\). Frobenius at \(y\) cycles the \(e\) copies of the stalk at \(x\), and its eigenvalues are \(e\)th roots of the eigenvalues upstairs. Measured against the residue field, absolute values are therefore unchanged (Lemma \hyperref[bp5]{18.1}).

The second fact is geometric, since restricting to a curve can lose monodromy. On \(\mathbf G_m\times\mathbf G_m\) with coordinates \(x,y\), in odd characteristic, the double cover \(u^2=x\) is connected, but over the curve \(x=1\) it splits into two copies. Accordingly, the Kummer sheaf of the quadratic character in \(x\) is geometrically nontrivial on the surface but constant on that curve. In the same way an irreducible sheaf of higher rank can become reducible on a badly chosen curve. On the other hand, a representation stays absolutely irreducible on a curve if the image of the fundamental group of the curve is large enough, and we show that this can be checked in one finite quotient (Lemma \hyperref[bp8]{18.2}). It then suffices that the curve stays connected in one finite étale cover, and we find such curves with Bertini's irreducibility theorem (Lemma \hyperref[bp11]{18.3}).

% object: lemma-finite-quasi-finite-direct-images (blueprint/chap-The-Deligne-Laumon-theorem-in-general/sec-Quasi-finite-morphisms-and-curves/lemma-finite-quasi-finite-direct-images/_section.md)
\phantomsection\label{bp5}%
\begin{bpstatementit}{Lemma~18.1}{ (Finite and quasi-finite direct images)}
Let \(g:Z_0\to Y_0\) be separated and quasi-finite. Then \(Rg_!\) is exact on sheaves and carries mixed sheaves of weights at most \(w\) to mixed sheaves of weights at most \(w\). It also carries \(D_{\leq w}(Z_0)\) into \(D_{\leq w}(Y_0)\). For a finite map, direct image carries a pure sheaf of weight \(w\) to a pure sheaf of weight \(w\).
\end{bpstatementit}

% proof: proof/lemma-finite-quasi-finite-direct-images
\begin{proof}[{Proof}]\phantomsection\label{bp6}%
We first treat a finite map and then reduce to it by Zariski's main theorem, which writes a quasi-finite map as an open immersion followed by a finite map. All the computation is in the finite case, since extension by zero does not change any nonzero stalk.

For a finite map, exactness follows from Lemma \bpout{D.1} and the comparison of coefficients. The same lemma shows that the stalk at a geometric point \(\bar y\) over a closed point \(y\) is the direct sum of the stalks at the geometric points of \(Z_0\) over \(\bar y\), and the local Frobenius at \(y\) permutes these summands. Let \(z\) be a point of residue degree \(e=[\kappa(z):\kappa(y)]\). Frobenius cycles through the \(e\) summands over \(z\), and its \(e\)th power on each of them is conjugate to the local Frobenius at \(z\). Each eigenvalue \(\beta\) on this cycle therefore satisfies \(\beta^e=\alpha\) for an eigenvalue \(\alpha\) at \(z\). If the input has weight \(w\), then

\[
|\iota(\beta)|^e=|\iota(\alpha)|=q_z^{w/2}=q_y^{ew/2},
\]

and we find \(|\iota(\beta)|=q_y^{w/2}\). This proves that finite direct image preserves purity. For example, for the constant sheaf on \(\operatorname{Spec}\mathbf F_{q^2}\) over \(\operatorname{Spec}\mathbf F_q\), the direct image has rank two and Frobenius swaps the two summands. Its eigenvalues \(\pm1\) have weight \(0\), as the constant sheaf does.

In general, \(Y_0\) is Noetherian, hence quasi-compact and quasi-separated, and Zariski's main theorem in the form [Stacks, \href{https://stacks.math.columbia.edu/tag/05K0}{05K0}] applies. It factors \(g\) as an open immersion \(j:Z_0\hookrightarrow\overline Z_0\) followed by a finite map \(\bar g:\overline Z_0\to Y_0\). By Proposition \bpout{5.28} we have \(Rg_!=\bar g_*j_!\), and both functors are exact. Extension by zero preserves purity, since its other stalks are zero. Now take a mixed sheaf of weights at most \(w\) on \(Z_0\), with a finite filtration by Weil subsheaves whose quotients are pure of weights at most \(w\). Since \(j\) and \(\bar g\) are defined over \(k_0\), the functors \(j_!\) and \(\bar g_*\) take Weil subsheaves to Weil subsheaves. Being exact, they take the filtration to a finite filtration of the direct image whose quotients are the direct images of the pure quotients. By the first part these are pure of the same weights, which proves the assertion on sheaves. Since \(Rg_!\) is exact, it commutes with taking cohomology sheaves, and the assertion on \(D_{\leq w}\) follows.
\end{proof}

% text: prose-18-03
\phantomsection\label{bp7}%
A lisse sheaf stays geometrically irreducible on a curve \(C\) in \(X\) as soon as the image of \(\pi_1(C)\) is large enough. The next lemma shows that it suffices for this image to surject onto one finite quotient of \(\pi_1(X)\), and in Lemma \hyperref[bp11]{18.3} we construct curves with this property.

% object: lemma-detecting-irreducibility-finite-quotient (blueprint/chap-The-Deligne-Laumon-theorem-in-general/sec-Quasi-finite-morphisms-and-curves/lemma-detecting-irreducibility-finite-quotient/_section.md)
\phantomsection\label{bp8}%
\begin{bpstatementit}{Lemma~18.2}{ (Detecting irreducibility in a finite quotient)}
Let \(\Pi\) be profinite and let \(\rho:\Pi\to\GL_E(V)\) be continuous and absolutely irreducible, with \(\dim_E V=n\). There are a \(\Pi\)-stable lattice \(L\subset V\) and an integer \(N\) such that every subgroup \(\Pi'\subseteq\Pi\) with the same image as \(\Pi\) in \(\GL(L/\pi_E^NL)\) acts absolutely irreducibly on \(V\).
\end{bpstatementit}

% proof: proof/lemma-detecting-irreducibility-finite-quotient
\begin{proof}[{Proof}]\phantomsection\label{bp9}%
Recall Burnside's theorem: an algebra of endomorphisms of a finite-dimensional vector space over an algebraically closed field that acts irreducibly is the full endomorphism algebra. Absolute irreducibility thus means that the group spans the full matrix algebra. We detect spanning by a nonzero determinant of \(n^2\) group elements. Such a determinant stays nonzero when we change the elements by a congruence modulo a high power of \(\pi_E\).

The compact group \(\rho(\Pi)\) stabilizes a lattice \(L\). Indeed, the stabilizer of any lattice \(L_0\) is open, hence of finite index in \(\rho(\Pi)\). So \(L_0\) has only finitely many translates, and their sum is a \(\rho(\Pi)\)-stable lattice. By absolute irreducibility and Burnside's theorem, there are elements \(g_1,\ldots,g_{n^2}\) whose matrices span \(\operatorname{End}_E(V)\). In a basis of \(L\), the determinant \(\Delta\) of their matrices, written as column vectors, is nonzero and integral. Choose \(N>v_{\pi_E}(\Delta)\). Suppose that \(g'_1,\ldots,g'_{n^2}\) have matrices congruent to those of the \(g_i\) modulo \(\pi_E^N\). Since the determinant is a polynomial with integer coefficients in the entries, their determinant \(\Delta'\) satisfies \(\Delta'\equiv\Delta\pmod{\pi_E^N}\). Hence \(v_{\pi_E}(\Delta')=v_{\pi_E}(\Delta)\), and in particular \(\Delta'\ne0\). A subgroup as in the statement contains such elements. Its matrices therefore span the whole matrix algebra, also after any field extension, and by Burnside's theorem it acts absolutely irreducibly. Finally, the image of \(\Pi\) in \(\GL(L/\pi_E^NL)\) is a finite quotient of \(\Pi\), since by continuity the kernel is open.
\end{proof}

% text: prose-18-04
\phantomsection\label{bp10}%
For a lisse sheaf on \(X\) the finite quotient corresponds to a connected finite Galois cover \(T\to X\). If \(C\subset X\) is a curve with \(T\times_XC\) connected, then the fundamental group of \(C\) surjects onto the Galois group of \(T\). What we need, then, are curves through given points whose inverse image in a given connected cover stays connected.

% object: lemma-curve-preserving-finite-cover (blueprint/chap-The-Deligne-Laumon-theorem-in-general/sec-Quasi-finite-morphisms-and-curves/lemma-curve-preserving-finite-cover/_section.md)
\phantomsection\label{bp11}%
\begin{bpstatementit}{Lemma~18.3}{ (A curve preserving a finite cover)}
Let \(X\) be a smooth connected affine scheme over \(\overline{\mathbf F}_q\), of dimension at least two. Let \(T\to X\) be a connected finite étale cover and let \(S\subset X\) be a nonempty finite set of closed geometric points. There is a smooth connected closed affine curve \(C\subset X\) containing \(S\), with \(T\times_XC\) connected. The curve and its inclusion are defined over a finite field.
\end{bpstatementit}

% proof: proof/lemma-curve-preserving-finite-cover
\begin{proof}[{Proof}]\phantomsection\label{bp12}%
We cut \(X\) by a general hypersurface through \(S\), show that it is smooth and that both it and its inverse image in \(T\) are connected, and then repeat.

\emph{The linear system.} Embed \(X\) in a projective closure \(\overline X\). For large \(n\), the sections of \(\mathcal O_{\overline X}(n)\) vanishing on \(S\) have base locus exactly \(S\) on \(X\) and define an immersion off \(S\) ([Stacks, \href{https://stacks.math.columbia.edu/tag/0FD5}{0FD5}]). They also surject onto the first-order jets \(\bigoplus_{s\in S}\mathfrak m_s/\mathfrak m_s^2\otimes
\mathcal O_{\overline X}(n)\) at \(S\). Indeed, if \(\mathcal I\) is the ideal sheaf of \(S\), then \(H^1(\overline X,\mathcal I^2(n))=0\) for large \(n\) by Serre vanishing ([Stacks, \href{https://stacks.math.columbia.edu/tag/0B5U}{0B5U}]), and the cohomology sequence of \(0\to\mathcal I^2(n)\to\mathcal I(n)\to(\mathcal I/\mathcal I^2)(n)\to0\) shows that \(H^0(\overline X,\mathcal I(n))\) surjects onto the jets.

\emph{Smoothness.} A general divisor of these sections is smooth on \(X\). The immersion off \(S\) says nothing about the points of \(S\) themselves, and there smoothness is controlled by the jets. A divisor through \(s\) is smooth at \(s\) exactly when its equation has nonzero image in \(\mathfrak m_s/\mathfrak m_s^2\). Since the sections surject onto the jets, the condition that a section have nonzero derivative at each \(s\) is nonempty and open. On \(X\setminus S\) the divisors are the pullbacks of hyperplanes under the immersion of \(X\setminus S\) into a projective space that these sections define. By Bertini's theorem ([Stacks, \href{https://stacks.math.columbia.edu/tag/0FD6}{0FD6}]) a general one is smooth there.

\emph{Connectedness.} A general divisor and its inverse image in \(T\) are irreducible. We use Bertini's irreducibility theorem [Stacks, \href{https://stacks.math.columbia.edu/tag/0G4F}{0G4F}]: if \(Z\) is an irreducible scheme of finite type over an algebraically closed field and \(Z\to\mathbf P^N\) is a morphism whose image has closure of dimension at least two, then the inverse image of a general hyperplane is irreducible. We apply it to the immersion of \(X\setminus S\) above and to its composition with the restriction of \(T\to X\) over \(X\setminus S\). In both cases the hyperplanes correspond to the sections vanishing on \(S\) and their pullbacks to \(T\). Both sources are smooth and connected, hence irreducible. Both images have closure of dimension \(\dim X\ge2\), because \(T\to X\) is finite and surjective. So for a general divisor \(D\), both \(D\cap(X\setminus S)\) and its inverse image in \(T\) are irreducible. The finitely many points over \(S\) cannot be components, since every component of a divisor has dimension \(\dim X-1\ge1\), and the claim follows.

\emph{Induction.} Intersecting these nonempty open conditions in the space of sections, we obtain a smooth connected affine hypersurface \(H\) containing \(S\) whose inverse image in \(T\) is connected. It is smooth, connected and affine of dimension \(\dim X-1\), it contains \(S\), and \(T\times_XH\to H\) is a connected finite étale cover. Thus \(H\) again satisfies the hypotheses of the lemma, and we repeat the construction until the dimension is one. The curve is defined over a finite field, because it has finitely many equations, whose coefficients lie in a finite extension of \(\mathbf F_q\).
\end{proof}

% text: prose-18-05
\phantomsection\label{bp13}%
Combining the two lemmas, we obtain the curves we need for real sheaves. We can treat several sheaves at once, because a single cover dominates the finite quotients attached to all of them.

% object: corollary-curve-preserving-geometric-irreducibility (blueprint/chap-The-Deligne-Laumon-theorem-in-general/sec-Quasi-finite-morphisms-and-curves/corollary-curve-preserving-geometric-irreducibility/_section.md)
\phantomsection\label{bp14}%
\begin{bpstatementit}{Corollary~18.4}{ (A curve preserving geometric irreducibility)}
Let \(X_0\) be a smooth geometrically connected affine scheme of positive dimension. Given finitely many geometrically irreducible lisse Weil sheaves on \(X_0\) and finitely many closed geometric points, there is, after a finite extension of the constant field, a smooth geometrically connected affine curve through those points on which all the sheaves remain geometrically irreducible.
\end{bpstatementit}

% proof: proof/corollary-curve-preserving-geometric-irreducibility
\begin{proof}[{Proof}]\phantomsection\label{bp15}%
We translate irreducibility into the surjectivity of \(\pi_1(C)\) onto one finite group, and then find \(C\) with Bertini's theorem applied to the corresponding cover.

Choose a finite extension \(E\) of \(\mathbf Q_\ell\) over which all the sheaves are defined, and apply Lemma \hyperref[bp8]{18.2} to their geometric representations. Let \(T\to X\) be the connected Galois cover corresponding to the intersection of the kernels of the resulting finite quotients of \(\pi_1(X)\), and let \(Q\) be its group. In dimension one, we take \(C=X\). Otherwise \(\dim X\ge2\), and we apply Lemma \hyperref[bp11]{18.3} to \(T\) and the given points. If no points are given we add one, since the lemma needs \(S\) nonempty. As \(T\times_XC\) is connected, \(\pi_1(C)\) acts transitively on the fibre of the \(Q\)-torsor \(T\times_XC\to C\), so its image in \(Q\) is all of \(Q\). Each of the finite quotients above is a quotient of \(Q\). Hence \(\pi_1(C)\) has the same image as \(\pi_1(X)\) in each of them, and Lemma \hyperref[bp8]{18.2}, applied with \(\Pi=\pi_1(X)\) and \(\Pi'\) the image of \(\pi_1(C)\), shows that the restrictions remain geometrically irreducible. By Lemma \hyperref[bp11]{18.3} the curve and the points are defined over a finite extension of the constant field. The curve \(C\) is geometrically connected because it is connected over \(\overline{\mathbf F}_q\).
\end{proof}


% object: sec-Mixedness-of-real-sheaves (blueprint/chap-The-Deligne-Laumon-theorem-in-general/sec-Mixedness-of-real-sheaves/_section.md), content hidden
\subsection*{\texorpdfstring{18.2\quad Mixedness of real sheaves}{18.2 Mixedness of real sheaves}}
\phantomsection\label{bp16}%
\addcontentsline{toc}{subsection}{\texorpdfstring{18.2\quad Mixedness of real sheaves}{18.2 Mixedness of real sheaves}}
% text: prose-18-06
\phantomsection\label{bp17}%
With these curves we pass from the purity theorem for real sheaves on curves (Theorem \bpout{14.26}) to the mixedness of real sheaves on any base. The curves preserve geometric irreducibility, so we first split a lisse sheaf into geometrically irreducible pieces. These pieces need not be Weil sheaves, because an irreducible Weil sheaf need not be geometrically irreducible. For example, let \(\chi\) be a character of \(\mathbf F_{q^2}^\times\) with \(\chi^q\ne\chi\), and let \(g\) be the map from \(\mathbf G_m\) over \(\mathbf F_{q^2}\) to \(\mathbf G_m\) over \(\mathbf F_q\). Geometrically, the direct image under \(g\) of the Kummer sheaf of \(\chi\) splits into the Kummer sheaves of \(\chi\) and \(\chi^q\), which are distinct. Frobenius exchanges them, and so the direct image is an irreducible Weil sheaf of rank two. After the quadratic constant extension the two pieces are Weil sheaves themselves. In general, the geometric constituents become Weil subquotients after a finite constant extension (Lemma \hyperref[bp18]{18.5}). In the proof of Theorem \hyperref[bp25]{18.8} we descend mixedness back along this extension by a trace retraction. We also reduce constructible sheaves to lisse sheaves on smooth affine strata, and glue the strata using Lemma \hyperref[bp5]{18.1}.

% object: lemma-geometric-constituents-after-constant-extension (blueprint/chap-The-Deligne-Laumon-theorem-in-general/sec-Mixedness-of-real-sheaves/lemma-geometric-constituents-after-constant-extension/_section.md)
\phantomsection\label{bp18}%
\begin{bpstatementit}{Lemma~18.5}{ (Geometric constituents after a constant extension)}
A lisse Weil sheaf on a smooth geometrically connected scheme admits, after one finite constant extension, a finite Weil filtration with geometrically irreducible quotients.
\end{bpstatementit}

% proof: proof/lemma-geometric-constituents-after-constant-extension
\begin{proof}[{Proof}]\phantomsection\label{bp19}%
An irreducible Weil representation is geometrically semisimple. A power of Frobenius fixes each isotypic component, and on it we can triangularize this power. The flags obtained in this way are stable under the Weil group of the constant extension, and by them we refine a composition series to the required filtration.

In detail, take a composition series for the Weil representation. By Lemma \bpout{14.11}, applied to the geometric fundamental group, which is normal in the Weil group, the geometric restriction of each irreducible factor is semisimple. Frobenius permutes the finitely many geometric isotypic components of the factor, and we choose a lift \(s\) of geometric Frobenius to the Weil group and a power \(s^m\) fixing all of them. On a component \(S\otimes M\), geometric monodromy acts irreducibly on \(S\) and trivially on \(M\). Since \(s^m\) fixes the isomorphism type of \(S\), there is an isomorphism \(u\) from \(S\) to its conjugate by \(s^m\). We claim that the action of \(s^m\) on the component has the form \(u\otimes v\) with \(v\) invertible on \(M\). Indeed, the composite of the action of \(s^m\) with \(u^{-1}\otimes1\) commutes with geometric monodromy, and by Schur's lemma the endomorphisms of \(S\otimes M\) commuting with it are \(1\otimes\operatorname{End}(M)\). Triangularizing \(v\), we obtain a stable flag whose geometric factors are \(S\).

We take a common \(m\) for the finitely many factors. Then the subgroup generated by geometric monodromy and \(s^m\) is the Weil group after the constant extension of degree \(m\). Choose a finite extension \(E'\) of the coefficient field over which these flags are defined. Since a subspace defined over \(E'\) is stable under this group if and only if its extension to \(\overline{\mathbf Q}_\ell\) is, the flags define lisse Weil subsheaves, and refining the composition series by them gives the required filtration.
\end{proof}

% text: prose-18-07
\phantomsection\label{bp20}%
In the example above, after the quadratic extension the sheaf \(g_*\mathcal K_\chi\) becomes the direct sum of the Kummer sheaves of \(\chi\) and \(\chi^q\), and either summand gives a filtration as in the lemma. Beyond this example, we use the lemma, together with the curves of Corollary \hyperref[bp14]{18.4}, to prove that real sheaves are mixed (Theorem \hyperref[bp25]{18.8}). In that proof, however, the filtrations with pure quotients exist only after a constant extension and only on the pieces of a stratification, so the next two lemmas recover mixedness of the original sheaf.

% object: lemma-mixedness-descends-constant-extension (blueprint/chap-The-Deligne-Laumon-theorem-in-general/sec-Mixedness-of-real-sheaves/lemma-mixedness-descends-constant-extension/_section.md)
\phantomsection\label{bp21}%
\begin{bpstatementit}{Lemma~18.6}{ (Mixedness and finite constant extensions)}
Let \(X_0\) be a separated scheme of finite type over \(k_0\), let \(g:X'_0=X_0\otimes_{k_0}k_m\to X_0\) be a finite constant extension, and let \(G\) be a constructible Weil sheaf on \(X_0\). If \(g^*G\) is mixed, then \(G\) is mixed.
\end{bpstatementit}

% proof: proof/lemma-mixedness-descends-constant-extension
\begin{proof}[{Proof}]\phantomsection\label{bp22}%
Since \(g^*G\) is mixed, so is \(g_*g^*G\) by Lemma \hyperref[bp5]{18.1}. The unit \(G\to g_*g^*G\) has the Weil-equivariant retraction \((\deg g)^{-1}\operatorname{Tr}\). Indeed, étale locally the unit is the diagonal and the trace is the sum over the copies. These local maps glue because the sum is invariant under permutations of the copies. So \(G\) is a Weil subsheaf of \(g_*g^*G\). A Weil subobject of a mixed sheaf is mixed (Proposition \bpout{13.7}), and therefore \(G\) is mixed.
\end{proof}

% object: lemma-mixedness-open-closed (blueprint/chap-The-Deligne-Laumon-theorem-in-general/sec-Mixedness-of-real-sheaves/lemma-mixedness-open-closed/_section.md)
\phantomsection\label{bp23}%
\begin{bpstatementit}{Lemma~18.7}{ (Mixedness on an open and its complement)}
Let \(j:U_0\to X_0\) be an open immersion with closed complement \(i:Z_0\to X_0\), and let \(G\) be a constructible Weil sheaf on \(X_0\). If \(j^*G\) and \(i^*G\) are mixed, then \(G\) is mixed.
\end{bpstatementit}

% proof: proof/lemma-mixedness-open-closed
\begin{proof}[{Proof}]\phantomsection\label{bp24}%
The sheaf \(G\) is an extension

\[
0\longrightarrow j_!j^*G\longrightarrow G
\longrightarrow i_*i^*G\longrightarrow0,
\]

whose outer terms are mixed by Lemma \hyperref[bp5]{18.1}, since \(j\) and \(i\) are quasi-finite. An extension of mixed sheaves is mixed (Proposition \bpout{13.7}).
\end{proof}

% object: theorem-real-constructible-sheaves-are-mixed (blueprint/chap-The-Deligne-Laumon-theorem-in-general/sec-Mixedness-of-real-sheaves/theorem-real-constructible-sheaves-are-mixed/_section.md)
\phantomsection\label{bp25}%
\begin{bpstatementit}{Theorem~18.8}{ (Real constructible sheaves are mixed)}
Every constructible \(\iota\)-real Weil sheaf on a separated finite-type \(k_0\)-scheme is \(\iota\)-mixed.
\end{bpstatementit}

% proof: proof/theorem-real-constructible-sheaves-are-mixed
\begin{proof}[{Proof}]\phantomsection\label{bp26}%
We join a fixed point \(x_*\) to an arbitrary point \(x\) by a curve on which the geometrically irreducible constituents stay irreducible. On that curve each constituent is pure by Theorem \bpout{14.26}. Its weight can be read off at \(x_*\), so it does not depend on the curve. This handles a lisse sheaf on a smooth affine scheme, and stratifying handles an arbitrary constructible sheaf. Descent along constant extensions and gluing of strata preserve mixedness (Lemmas \hyperref[bp21]{18.6} and \hyperref[bp23]{Lemma 18.7}), so nothing about curves is used beyond Theorem \bpout{14.26}.

\emph{The lisse case.} Suppose that \(X_0\) is smooth, geometrically connected and affine, and that the sheaf \(G\) is lisse. Since the curves of Corollary \hyperref[bp14]{18.4} preserve geometric irreducibility, we first split \(G\) into geometrically irreducible pieces. To do so, we make the constant extension of Lemma \hyperref[bp18]{18.5} and write \(S_1,\ldots,S_r\) for the quotients of its Weil filtration. Note that the sheaf \(G\) remains real after this extension, because the powers of a conjugation-stable multiset of eigenvalues again form a conjugation-stable multiset.

Fix a closed point \(x_*\) of \(X_0\), and let \(x\) be any other closed point. By Corollary \hyperref[bp14]{18.4}, there is a curve through chosen geometric points above \(x_*\) and \(x\) on which every \(S_j\) remains irreducible. The restriction of \(G\) to this curve is real, and Theorem \bpout{14.26} says that the irreducible constituents of a real sheaf on a curve are pure. Since the restriction of \(S_j\) is irreducible, it is one of them, and so it is pure. Let \(\alpha\) be an eigenvalue of \(S_j\) at \(x_*\), and let \(e\) be the residue degree over \(\kappa(x_*)\) of the point of the curve above \(x_*\). At that point, then, the corresponding eigenvalue is \(\alpha^e\) and the residue cardinality is \(q_{x_*}^e\). If the restriction is pure of weight \(u\), then \(|\iota(\alpha)|^e=q_{x_*}^{eu/2}\), so the weight of the restriction of \(S_j\) must be

\[
w_j=2\log_{q_{x_*}}|\iota(\alpha)|,
\]

which does not depend on the curve. In particular, all eigenvalues of \(S_j\) at \(x_*\) have this modulus. At \(x\), the same calculation with powers shows that \(S_j\) is pure of weight \(w_j\). As \(x\) was arbitrary, each \(S_j\) is pure of weight \(w_j\) on \(X_0\). The chosen Weil filtration then shows that \(G\) is mixed after the constant extension, and by Lemma \hyperref[bp21]{18.6}, \(G\) is mixed.

\emph{The general case.} Now let \(G\) be an arbitrary constructible real sheaf. Since nilpotents do not change the étale site, we may assume that \(X_0\) is reduced. In each irreducible component, choose a dense smooth affine open on which \(G\) is lisse and which avoids the other components. Since the complement has smaller dimension, iterating we obtain a finite chain of closed subschemes whose differences are disjoint unions of smooth connected affine pieces on which \(G\) is lisse. After one finite constant extension these pieces are geometrically connected, and such an extension is allowed by Lemma \hyperref[bp21]{18.6}, and it keeps \(G\) real. The lisse case then applies to the positive-dimensional pieces, while at a zero-dimensional piece a Frobenius-stable full flag of each stalk has quotients of rank one, which are pure. The restriction of \(G\) to every piece is therefore mixed, and Lemma \hyperref[bp23]{18.7}, applied successively along the chain of closed subschemes, shows that \(G\) is mixed.
\end{proof}


% object: sec-Relative-affine-curves (blueprint/chap-The-Deligne-Laumon-theorem-in-general/sec-Relative-affine-curves/_section.md), content hidden
\subsection*{\texorpdfstring{18.3\quad Relative affine curves}{18.3 Relative affine curves}}
\phantomsection\label{bp27}%
\addcontentsline{toc}{subsection}{\texorpdfstring{18.3\quad Relative affine curves}{18.3 Relative affine curves}}
% text: prose-18-08
\phantomsection\label{bp28}%
Let \(h:U_0\to Y_0\) be a smooth family of affine curves and \(F\) a pure lisse sheaf on \(U_0\). On each stalk of \(R^ih_!F\), Theorem \bpout{17.15} bounds the eigenvalues by \(w+i\). Suppose that \(R^ih_!F\) is mixed. Each pure quotient of its filtration has a nonzero stalk somewhere, and its eigenvalues there are among those of \(R^ih_!F\). By Lemma \bpout{13.3} the quotients therefore have weights at most \(w+i\). Most of the proof of Theorem \hyperref[bp29]{18.9} is devoted to showing that \(R^ih_!F\) is mixed.

% object: theorem-pure-lisse-sheaves-over-relative (blueprint/chap-The-Deligne-Laumon-theorem-in-general/sec-Relative-affine-curves/theorem-pure-lisse-sheaves-over-relative/_section.md)
\phantomsection\label{bp29}%
\begin{bpstatementit}{Theorem~18.9}{ (Pure lisse sheaves over relative affine curves)}
Let \(h:U_0\to Y_0\) be smooth separated of relative dimension one, with affine geometric fibres, and let \(F\) be a lisse Weil sheaf on \(U_0\), pure of weight \(w\). Then \(R^ih_!F\) is mixed of weights at most \(w+i\).
\end{bpstatementit}

% proof: proof/theorem-pure-lisse-sheaves-over-relative
\begin{proof}[{Proof}]\phantomsection\label{bp30}%
We bound the stalks of \(R^ih_!F\) by the curve case, we show that these sheaves are mixed by twisting to weight zero and adding the dual, and we combine the two facts using Lemma \bpout{13.3}.

\emph{The stalks.} Fix a closed point \(y\) of \(Y_0\). Base change for direct images with compact support (Proposition \bpout{5.28}) gives an identification

\[
(R^ih_!F)_{\bar y}=H_c^i(U_{\bar y},F)
\]

compatible with the local Frobenius operators. The restriction of \(F\) to the fibre is again pure of weight \(w\), because a residue extension raises both an eigenvalue and the residue cardinality to the same power. The fibre is a smooth affine curve over \(\kappa(y)\), so Theorem \bpout{17.15} shows that \(H_c^i(U_{\bar y},F)\) has weights at most \(w+i\). After a finite extension of \(\kappa(y)\) the components of the fibre are geometrically connected. Lemma \bpout{15.7}, applied to each of them with \(j\) the identity, shows that these groups vanish outside degrees \(1\) and \(2\).

\emph{Mixedness.} We prove that \(R^ih_!F\) is mixed. This is trivial for \(F=0\), so we may assume \(F\ne0\). Let \(G=F^0\oplus(F^0)^\vee\), with \(F^0=F\otimes\mathcal B_b\), be the sheaf of Proposition \bpout{15.9}. It is lisse, real and pure of weight zero. At a closed point \(y\) write

\[
P_{i,y}(T)=\det(1-T\operatorname{Frob}_y\mid H_c^i(U_{\bar y},G)).
\]

Applying Lemma \bpout{15.7} to the components of the fibre after a finite extension of \(\kappa(y)\), as above, we see that \(P_{0,y}=1\) and that every reciprocal root of \(P_{2,y}\) has modulus \(q_y\). By the bound just proved, every reciprocal root of \(P_{1,y}\) has modulus at most \(q_y^{1/2}\).

We claim that \(\iota P_{1,y}\) and \(\iota P_{2,y}\) have real coefficients. By the Euler product and the determinant identity of Theorem \bpout{12.30}, after applying \(\iota\) we have

\[
L_\iota(U_y,G,T)=\frac{\iota P_{1,y}(T)}{\iota P_{2,y}(T)}.
\]

Since the Euler factors are real, this rational function has real Taylor coefficients. Both it and its complex conjugate are regular at zero, so they are equal as rational functions. Moreover, the numerator and the denominator are coprime: a common root would give an eigenvalue having both permitted moduli, whereas \(q_y^{1/2}<q_y\). Since a reduced fraction whose numerator and denominator have constant coefficient one is unique, comparing it with its conjugate we find \(\iota P_{1,y},\iota P_{2,y}\in\mathbf R[T]\), which proves the claim. Without the curve case this step would fail, because a common factor of numerator and denominator could cancel, and we would only know that the quotient is real, not each polynomial.

The sheaves \(R^ih_!G\) are constructible by the finiteness theorems, and we have shown that they are real at every closed point. By Theorem \hyperref[bp25]{18.8} each of them has a finite mixed filtration.

To return from \(G\) to \(F\), note that by additivity \(R^ih_!F^0\) is a direct summand of \(R^ih_!G\). A direct summand is a Weil subsheaf, and so it is mixed by Proposition \bpout{13.7}. By the projection formula for the constant twist, it is isomorphic to \(R^ih_!F\otimes\mathcal B_b\). Removing the twist preserves the finite Weil filtration and shifts the weight of every pure quotient by \(w\). Hence \(R^ih_!F\) is mixed.

\emph{Conclusion.} The sheaf \(R^ih_!F\) is mixed, and by the bound on the stalks all its local eigenvalues have modulus at most \(q_y^{(w+i)/2}\). Lemma \bpout{13.3}, with \(w+i\) in place of \(w\), shows that every nonzero quotient in the filtration has weight at most \(w+i\).
\end{proof}


% object: sec-Proof-of-the-Deligne-Laumon-theorem (blueprint/chap-The-Deligne-Laumon-theorem-in-general/sec-Proof-of-the-Deligne-Laumon-theorem/_section.md), content hidden
\subsection*{\texorpdfstring{18.4\quad Proof of the Deligne–Laumon theorem}{18.4 Proof of the Deligne–Laumon theorem}}
\phantomsection\label{bp31}%
\addcontentsline{toc}{subsection}{\texorpdfstring{18.4\quad Proof of the Deligne–Laumon theorem}{18.4 Proof of the Deligne–Laumon theorem}}
% text: prose-18-09
\phantomsection\label{bp32}%
An arbitrary \(f\) reduces to relative affine curves and quasi-finite maps. Over a field, Noether normalization makes an affine variety of dimension \(d\) finite over \(\mathbf A^d\), and \(\mathbf A^d\to\operatorname{Spec}k\) is a composite of \(d\) projections of affine lines. For a separated morphism \(f\) of finite type the same works over a dense open of the target, and we treat the complement by induction on its dimension (Lemma \hyperref[bp42]{18.14}). So we must check the weight bound for two kinds of maps and show that it survives the gluing.

Quasi-finite maps were treated in Lemma \hyperref[bp5]{18.1}. For the projection \(\mathbf A^1_Y\to Y\) (Proposition \hyperref[bp38]{18.12}), consider a sheaf on the relative line. Over a dense open of the base it is lisse off a closed subscheme finite over the base. It is therefore glued from a relative affine curve and a quasi-finite piece. The bound is stable under composition by the Leray spectral sequence (Theorem \hyperref[bp45]{18.15}). It is stable under partitions of the source and of the target by the localization sequences below. These gluing steps use only that mixed sheaves of weights at most \(u\) are closed under subquotients and extensions.

% object: lemma-weights-exact-sequence (blueprint/chap-The-Deligne-Laumon-theorem-in-general/sec-Proof-of-the-Deligne-Laumon-theorem/lemma-weights-exact-sequence/_section.md)
\phantomsection\label{bp33}%
\begin{bpstatementit}{Lemma~18.10}{ (Weights in an exact sequence)}
Let \(F'\to F\to F''\) be an exact sequence of Weil sheaves. If \(F'\) and \(F''\) are mixed of weights at most \(u\), then so is \(F\).
\end{bpstatementit}

% proof: proof/lemma-weights-exact-sequence
\begin{proof}[{Proof}]\phantomsection\label{bp34}%
Let \(a:F'\to F\) and \(b:F\to F''\) be the maps. Exactness at \(F\) gives a short exact sequence

\[
0\to\operatorname{im}(a)\to F\to\operatorname{im}(b)\to0,
\]

where \(\operatorname{im}(a)\) is a quotient of \(F'\) and \(\operatorname{im}(b)\) a subobject of \(F''\). Since mixed sheaves of weights at most \(u\) are closed under subquotients by Lemma \bpout{13.6} and under extensions by Proposition \bpout{13.7}, \(F\) is mixed of weights at most \(u\).
\end{proof}

% text: prose-18-10
\phantomsection\label{bp35}%
By the next lemma, to prove a weight bound for \(R^if_!\) of a morphism \(f\), it suffices to prove it over the pieces of finite partitions of the source and of the target.

% object: lemma-weight-bounds-partitions (blueprint/chap-The-Deligne-Laumon-theorem-in-general/sec-Proof-of-the-Deligne-Laumon-theorem/lemma-weight-bounds-partitions/_section.md)
\phantomsection\label{bp36}%
\begin{bpstatementit}{Lemma~18.11}{ (Weight bounds and partitions)}
Let \(f:X_0\to Y_0\) be separated of finite type, and let \(K\) be a constructible Weil sheaf on \(X_0\), mixed of weights at most \(w\).

\begin{enumerate}
\item Let \(j:U_0\to X_0\) be an open immersion with closed complement \(i:Z_0\to X_0\). If \(R^q(fj)_!j^*K\) and \(R^q(fi)_!i^*K\) are mixed of weights at most \(w+q\) for every \(q\), then so is \(R^qf_!K\).
\item Let \(V_0\subset Y_0\) be open with closed complement \(W_0\), and let \(f_V\) and \(f_W\) be the base changes of \(f\) to \(V_0\) and \(W_0\). If \(R^qf_{V!}(K|_{f^{-1}V_0})\) and \(R^qf_{W!}(K|_{f^{-1}W_0})\) are mixed of weights at most \(w+q\) for every \(q\), then so is \(R^qf_!K\).
\end{enumerate}
\end{bpstatementit}

% proof: proof/lemma-weight-bounds-partitions
\begin{proof}[{Proof}]\phantomsection\label{bp37}%
In both cases the restrictions of \(K\) are again mixed of weights at most \(w\) (Proposition \bpout{13.7}), and we place \(R^qf_!K\) in the middle of an exact sequence of Weil sheaves whose outer terms are bounded.

For (1), applying \(Rf_!\) and the composition isomorphism to the triangle \(j_!j^*K\to K\to i_*i^*K\to\) gives an exact sequence

\[
R^q(fj)_!j^*K\longrightarrow R^qf_!K\longrightarrow R^q(fi)_!i^*K,
\]

which is Frobenius-equivariant by naturality, and Lemma \hyperref[bp33]{18.10} gives the bound for \(R^qf_!K\).

For (2), let \(j':V_0\to Y_0\) and \(i':W_0\to Y_0\) be the inclusions. Base change for direct images with compact support (Proposition \bpout{5.28}) identifies \(j'^*R^qf_!K\) with \(R^qf_{V!}(K|_{f^{-1}V_0})\) and \(i'^*R^qf_!K\) with \(R^qf_{W!}(K|_{f^{-1}W_0})\). By Lemma \hyperref[bp5]{18.1}, the outer terms of

\[
0\longrightarrow j'_!j'^*R^qf_!K\longrightarrow R^qf_!K
\longrightarrow i'_*i'^*R^qf_!K\longrightarrow0
\]

are mixed of weights at most \(w+q\), and Lemma \hyperref[bp33]{18.10} gives the bound for the middle term.
\end{proof}

% object: proposition-relative-affine-line (blueprint/chap-The-Deligne-Laumon-theorem-in-general/sec-Proof-of-the-Deligne-Laumon-theorem/proposition-relative-affine-line/_section.md)
\phantomsection\label{bp38}%
\begin{bpstatementit}{Proposition~18.12}{ (A relative affine line)}
Let \(p:\mathbf A^1_{Y_0}\to Y_0\) be the projection, and let \(G\) be a constructible Weil sheaf on \(\mathbf A^1_{Y_0}\), mixed of weights at most \(w\). Then \(R^ip_!G\) is mixed of weights at most \(w+i\) for every \(i\).
\end{bpstatementit}

% proof: proof/proposition-relative-affine-line
\begin{proof}[{Proof}]\phantomsection\label{bp39}%
We stratify the target so that over each stratum, \(G\) is lisse off a subscheme that is finite over the base. On the lisse part the projection is a relative affine curve, and we apply Theorem \hyperref[bp29]{18.9}. On the finite part it is quasi-finite, and we apply Lemma \hyperref[bp5]{18.1}.

Take a finite mixed filtration of \(G\). In the long exact sequences of \(R^ip_!\) for its steps, the bound \(w+i\) passes from the pure quotients to \(G\) by Lemma \hyperref[bp33]{18.10}. We may therefore assume that \(G\) is pure and constructible of weight \(v\leq w\).

Suppose first that the target \(Y_0=\operatorname{Spec}R\) is integral and affine, so that \(\mathbf A^1_{Y_0}\) is integral. We construct the strata as follows. Choose a dense open \(U\) on which \(G\) is lisse. Since the locus where \(G\) is lisse is stable under Frobenius, we may replace \(U\) by the intersection of its finitely many Frobenius conjugates. Then \(U\) is defined over \(k_0\), and so are the ideal \(I\subset R[T]\) of its reduced complement and the choices below. If the complement is nonempty, choose a nonzero \(g\in I\) and invert its nonzero leading coefficient \(b\). Then over \(R_b\) the polynomial \(b^{-1}g\) is monic and vanishes in \(R_b[T]/I\). Hence the image of \(T\) is integral over \(R_b\), and \(R_b[T]/I\) is generated as an \(R_b\)-module by the powers \(T^k\) with \(k<\deg g\). In particular \(R_b[T]/I\) is finite over \(R_b\); if \(g\) has degree zero, the complement is empty over \(D(b)\).

Over \(D(b)\) the complement is therefore finite, hence quasi-finite. The open \(U\) is smooth of relative dimension one over \(D(b)\), and each of its fibres is an affine line minus finitely many points. So \(U\to D(b)\) satisfies the hypotheses of Theorem \hyperref[bp29]{18.9}. Applying that theorem on \(U\) and Lemma \hyperref[bp5]{18.1} on the complement, we obtain the bound over \(D(b)\) from Lemma lemma-weight-bounds-partitions.

For a general target we argue by induction on \(\dim Y_0\). Since nilpotents do not change the étale site, we may replace \(Y_0\) by its reduction. Take disjoint integral affine opens containing the generic points of its components, and perform this construction on each of them. Their union is open, and its complement has smaller dimension. The case of dimension zero is contained in the construction, because a reduced zero-dimensional scheme is a finite disjoint union of integral affine pieces. In general, the restriction of \(G\) over the closed complement is again mixed of weights at most \(w\), and the inductive hypothesis gives the bound there. Lemma lemma-weight-bounds-partitions, applied to the open union and its closed complement in the target, gives the bound over all of \(Y_0\).
\end{proof}

% object: example-leading-coefficient-stratification (blueprint/chap-The-Deligne-Laumon-theorem-in-general/sec-Proof-of-the-Deligne-Laumon-theorem/example-leading-coefficient-stratification/_section.md)
\phantomsection\label{bp40}%
\begin{bpstatement}{Example~18.13}{ (Why the leading coefficient is inverted)}
Let \(Y_0=\mathbf A^1_{k_0}=\operatorname{Spec}k_0[y]\), and suppose that \(G\) is lisse on the relative line \(\mathbf A^1_{Y_0}=\operatorname{Spec}k_0[y,T]\) off the curve \(Z=V(yT-1)\). This curve is not finite over the base. Indeed, over a point \(y\ne0\) the fibre of \(Z\) is the single point \(T=1/y\), and over \(y=0\) it is empty, since the equation becomes \(-1=0\). So \(Z\to Y_0\) is quasi-finite, but its image is \(D(y)\), which is not closed, whereas a finite map is closed; geometrically, as \(y\) tends to \(0\) the point \(1/y\) runs off to infinity. To repair this, the construction in the proof takes \(g=yT-1\) with leading coefficient \(b=y\). Over \(D(y)\) the polynomial \(T-y^{-1}\) is monic, and indeed \(k_0[y,y^{-1},T]/(yT-1)=k_0[y,y^{-1}]\) is finite over the base. The closed stratum \(y=0\), which has smaller dimension, is treated by induction, and on it the complement of the lisse locus happens to be empty.
\end{bpstatement}

% text: prose-18-11
\phantomsection\label{bp41}%
To apply the proposition to an arbitrary morphism, we factor it, over strata, as a quasi-finite map followed by projections of relative affine lines.

% object: lemma-relative-noether-normalization-over-stratification (blueprint/chap-The-Deligne-Laumon-theorem-in-general/sec-Proof-of-the-Deligne-Laumon-theorem/lemma-relative-noether-normalization-over-stratification/_section.md)
\phantomsection\label{bp42}%
\begin{bpstatementit}{Lemma~18.14}{ (Relative Noether normalization over a stratification)}
Let \(f:X_0\to Y_0\) be a morphism of finite type. There are finite partitions of \(X_0\) and of \(Y_0\) into locally closed pieces such that each piece of \(X_0\) lies over a piece \(V\) of \(Y_0\), and its map to \(V\) factors through a finite morphism to \(\mathbf A^d_V\) for some \(d\).
\end{bpstatementit}

% proof: proof/lemma-relative-noether-normalization-over-stratification
\begin{proof}[{Proof}]\phantomsection\label{bp43}%
We spread out Noether normalization from the generic points of the target, and then argue by induction on the dimension of the complement. Spreading out is possible because a normalization over the fraction field involves finitely many equations with finitely many denominators. Recall Noether normalization over a field \(k'\) [Stacks, \href{https://stacks.math.columbia.edu/tag/00OY}{00OY}]: a nonzero finitely generated \(k'\)-algebra is finite over a polynomial subalgebra.

Let \(R\) be an integral domain with fraction field \(k'\), let \(\operatorname{Spec}R\) be an affine open of the target, and let \(B\) be a finitely generated \(R\)-algebra. If \(B\otimes_Rk'=0\), then \(B_s=0\) for some nonzero \(s\in R\). Otherwise, Noether normalization over \(k'\) provides elements \(t_1,\ldots,t_d\) and monic equations for each of a fixed set of algebra generators of \(B\) over \(k'[t_1,\ldots,t_d]\). Choose a nonzero \(s\in R\) divisible by all denominators of the \(t_i\) and of the coefficients of these equations. Each equation, read in \(B_s\), has a left-hand side that vanishes in \(B\otimes_Rk'\) and is therefore killed by a nonzero element of \(R\). Multiplying \(s\) by these finitely many elements, we may assume that the equations hold in \(B_s\). Then bounded powers of the generators span \(B_s\) over \(R_s[t_1,\ldots,t_d]\), and the resulting map to \(\mathbf A^d_{R_s}\) is finite.

To obtain the partitions, choose disjoint integral affine opens containing the generic points of the reduced target. Partition their inverse images into finitely many reduced affine locally closed pieces and apply the preceding construction to each piece. On each target open, a single product of the finitely many denominators then works for all pieces. Since the complement of the resulting open has smaller dimension, by induction on this complement we obtain finite partitions of the whole source and target. Moreover, all choices are made over \(k_0\).
\end{proof}

% text: prose-18-12
\phantomsection\label{bp44}%
On each piece the map to \(V\) is thus a finite map followed by \(d\) projections of relative affine lines. The bound holds for the finite map by Lemma \hyperref[bp5]{18.1} and for each projection by Proposition \hyperref[bp38]{18.12}, and the next theorem puts the pieces together.

% object: theorem-direct-image-mixed-sheaf (blueprint/chap-The-Deligne-Laumon-theorem-in-general/sec-Proof-of-the-Deligne-Laumon-theorem/theorem-direct-image-mixed-sheaf/_section.md)
\phantomsection\label{bp45}%
\begin{bpstatementit}{Theorem~18.15}{ (Direct image of a mixed sheaf)}
Let \(f:X_0\to Y_0\) be separated of finite type and let \(F\) be mixed of weights at most \(w\). Then \(R^if_!F\) is mixed of weights at most \(w+i\) for every \(i\).
\end{bpstatementit}

% proof: proof/theorem-direct-image-mixed-sheaf
\begin{proof}[{Proof}]\phantomsection\label{bp46}%
Say that a morphism has the property if the conclusion holds for it, for every \(w\) and every mixed \(F\) of weights at most \(w\). Quasi-finite maps and projections of relative affine lines have the property, and by Lemma \hyperref[bp42]{18.14} every \(f\) is built from these two kinds of maps. So it suffices to show that the property is stable under composition and under finite partitions.

For composition, let \(X_0\xrightarrow{h}Z_0\xrightarrow{g}Y_0\) and consider the Leray spectral sequence for compact supports,

\[
E_2^{a,b}=R^ag_!(R^bh_!F)\ \Longrightarrow\ R^{a+b}(gh)_!F.
\]

It is Frobenius-equivariant by the naturality of the truncation tower and of the composition isomorphism for direct image with compact support. If both maps have the property, then \(E_2^{a,b}\) is mixed of weights at most \(w+a+b\). Since each later page is a subquotient of the previous one, the same holds for \(E_\infty^{a,b}\). The abutment filtration of \(R^n(gh)_!F\) is finite, because the operations preserve bounded constructible complexes. Its quotients are the \(E_\infty^{a,n-a}\), and by Lemma \hyperref[bp33]{18.10} \(R^n(gh)_!F\) is mixed of weights at most \(w+n\).

For partitions, we apply Lemma \hyperref[bp42]{18.14} and refine the partitions to filtrations by open and closed pieces. After this refinement we can order the pieces so that each is open in the union of itself and the later ones. This lets us apply the localization sequences one piece at a time. On the pieces, the restriction of \(F\) to each source piece is again mixed of weights at most \(w\), and on every refined source piece the finite map to \(\mathbf A^d_V\) restricts to a separated quasi-finite map, which has the property by Lemma \hyperref[bp5]{18.1}. Moreover, the projection \(\mathbf A^d_V\to V\) is a composition of \(d\) projections of relative affine lines. Each of them has the property by Proposition \hyperref[bp38]{18.12}, and by the composition argument the map from the piece to \(V\) has the property. Finally, we assemble the pieces one at a time in the chosen order. By Lemma lemma-weight-bounds-partitions on the source and Lemma lemma-weight-bounds-partitions on the target, the bound for the pieces implies the bound for \(R^if_!F\) on \(Y_0\).
\end{proof}

% text: prose-18-13
\phantomsection\label{bp47}%
For Theorem \bpout{13.11}, stated in Section \bpout{13.3}, the case of sheaves is Theorem \hyperref[bp45]{18.15}, and the spectral sequence of the truncation tower passes to complexes.

% proof: proof/theorem-weights-rf-shriek
\begin{proof}[{Proof of Theorem 13.11}]\phantomsection\label{bp48}%
Let \(K\in D_{\le w}(X_0)\), so that each \(\mathcal H^b(K)\) is mixed of weights at most \(w+b\). Applying \(Rf_!\) to the ordinary truncation tower of \(K\), which is bounded and functorial, we obtain the spectral sequence

\[
E_2^{a,b}=R^af_!\mathcal H^b(K)
\ \Longrightarrow\ \mathcal H^{a+b}(Rf_!K),
\]

By naturality and the Frobenius compatibility of \(Rf_!\), its differentials and its finite abutment filtration are compatible with the Weil structures. By Theorem \hyperref[bp45]{18.15}, \(E_2^{a,b}\) is mixed of weights at most \(w+a+b\), and so is its subquotient \(E_\infty^{a,b}\). The quotients of the filtration of \(\mathcal H^n(Rf_!K)\) are the \(E_\infty^{a,n-a}\). By Lemma \hyperref[bp33]{18.10} and induction on the length of the filtration, \(\mathcal H^n(Rf_!K)\) is mixed of weights at most \(w+n\). In other words, \(Rf_!K\in D_{\le w}(Y_0)\). The last assertion is the case of the structure map of \(X_0\), for which the geometric stalk of \(R^nf_!K\) is \(H_c^n(X,K)\) with its Frobenius operator.
\end{proof}



\clearpage\phantomsection\hypertarget{bp-outside}{}%
\section*{Outside this export}
\addcontentsline{toc}{section}{Outside this export}
References marked \ensuremath{{}^\ast} are to parts of the blueprint this document leaves out:
\begin{description}
\item[Proposition 5.28] Composition and base change for $Rf_!$, in Ⅰ The six operations (\bpslug{proposition-{}\allowbreak{}composition-{}\allowbreak{}base-{}\allowbreak{}change-{}\allowbreak{}rf-{}\allowbreak{}shriek})
\item[Theorem 12.30] The trace formula as an Euler product, in Ⅱ Duality and the trace formula (\bpslug{theorem-{}\allowbreak{}trace-{}\allowbreak{}formula-{}\allowbreak{}euler-{}\allowbreak{}product})
\item[Definition 13.1] Weights, purity and mixedness, in Ⅲ Weights and the Riemann hypothesis (\bpslug{definition-{}\allowbreak{}weights})
\item[Lemma 13.3] Bounding the weights of a mixed sheaf, in Ⅲ Weights and the Riemann hypothesis (\bpslug{lemma-{}\allowbreak{}bounding-{}\allowbreak{}weights-{}\allowbreak{}mixed-{}\allowbreak{}sheaf})
\item[Lemma 13.6] Subobjects and quotients of Weil sheaves, in Ⅲ Weights and the Riemann hypothesis (\bpslug{lemma-{}\allowbreak{}subobjects-{}\allowbreak{}quotients-{}\allowbreak{}weil-{}\allowbreak{}sheaves})
\item[Proposition 13.7] Weights under the usual operations, in Ⅲ Weights and the Riemann hypothesis (\bpslug{proposition-{}\allowbreak{}weights-{}\allowbreak{}under-{}\allowbreak{}usual-{}\allowbreak{}operations})
\item[Section 13.3] The Deligne–Laumon theorem, in Ⅲ Weights and the Riemann hypothesis (\bpslug{sec-{}\allowbreak{}The-{}\allowbreak{}Deligne-{}\allowbreak{}Laumon-{}\allowbreak{}theorem-{}\allowbreak{}2})
\item[Theorem 13.11] Weights of $Rf_!$, in Ⅲ Weights and the Riemann hypothesis (\bpslug{theorem-{}\allowbreak{}weights-{}\allowbreak{}rf-{}\allowbreak{}shriek})
\item[Lemma 14.11] Restriction to a normal subgroup, in Ⅲ Weights and the Riemann hypothesis (\bpslug{lemma-{}\allowbreak{}restriction-{}\allowbreak{}normal-{}\allowbreak{}subgroup})
\item[Theorem 14.26] Purity of real sheaves on curves, in Ⅲ Weights and the Riemann hypothesis (\bpslug{theorem-{}\allowbreak{}purity-{}\allowbreak{}real-{}\allowbreak{}sheaves-{}\allowbreak{}curves})
\item[Lemma 15.7] The extreme degrees on an affine curve, in Ⅲ Weights and the Riemann hypothesis (\bpslug{lemma-{}\allowbreak{}extreme-{}\allowbreak{}degrees-{}\allowbreak{}affine-{}\allowbreak{}curve})
\item[Proposition 15.9] Twisting to weight zero and adding the dual, in Ⅲ Weights and the Riemann hypothesis (\bpslug{proposition-{}\allowbreak{}twisting-{}\allowbreak{}weight-{}\allowbreak{}zero-{}\allowbreak{}adding-{}\allowbreak{}dual})
\item[Theorem 17.15] The weight bound on a curve, in Ⅲ Weights and the Riemann hypothesis (\bpslug{theorem-{}\allowbreak{}weight-{}\allowbreak{}bound-{}\allowbreak{}curve})
\item[Lemma D.1] Finite direct images, in Appendices (\bpslug{lemma-{}\allowbreak{}finite-{}\allowbreak{}direct-{}\allowbreak{}images})
\end{description}
\clearpage\phantomsection%
\section*{References}
\addcontentsline{toc}{section}{References}
\bpbibitem{BBD82}{bp50}{A. Beilinson, J. Bernstein, and P. Deligne. \emph{Faisceaux pervers}, volume 100 of \emph{Astérisque}. 1982.}

\bpbibitem{BS15}{bp51}{B. Bhatt and P. Scholze. The pro-étale topology for schemes. \emph{Astérisque}, 369:99–201, 2015. \href{https://people.mpim-bonn.mpg.de/scholze/proetale.pdf}{Author's manuscript}.}

\bpbibitem{Con07}{bp52}{B. Conrad. Deligne's notes on Nagata compactifications. \emph{Journal of the Ramanujan Mathematical Society}, 22:205–257, 2007. \href{https://math.stanford.edu/\~conrad/papers/nagatafinal.pdf}{Author's manuscript}.}

\bpbibitem{Del74}{bp49}{P. Deligne. La conjecture de Weil. I. \emph{Publ. Math. IHÉS}, 43:273–307, 1974.}

\bpbibitem{Del80}{bp53}{P. Deligne. La conjecture de Weil. II. \emph{Publ. Math. IHÉS}, 52:137–252, 1980.}

\bpbibitem{Den22}{bp54}{C. Deninger. There is no “Weil-”cohomology theory with real coefficients for arithmetic curves. \emph{Ann. Sc. Norm. Super. Pisa Cl. Sci.}, 2022. To appear. arXiv:\href{https://arxiv.org/abs/2204.02714}{2204.02714}.}

\bpbibitem{Deu41}{bp55}{M. Deuring. Die Typen der Multiplikatorenringe elliptischer Funktionenkörper. \emph{Abh. Math. Sem. Univ. Hamburg}, 14:197–272, 1941.}

\bpbibitem{Eke90}{bp56}{T. Ekedahl. On the adic formalism. In \emph{The Grothendieck Festschrift, Vol. II}, volume 87 of \emph{Progress in Mathematics}, pages 197–218. Birkhäuser, 1990.}

\bpbibitem{Has36}{bp57}{H. Hasse. Zur Theorie der abstrakten elliptischen Funktionenkörper. I, II, III. \emph{J. Reine Angew. Math.}, 175:55–62, 69–88, 193–208, 1936.}

\bpbibitem{Hum72}{bp58}{J. E. Humphreys. \emph{Introduction to Lie Algebras and Representation Theory}, volume 9 of \emph{Graduate Texts in Mathematics}. Springer, 1972.}

\bpbibitem{Ill94}{bp59}{L. Illusie. Autour du théorème de monodromie locale. \emph{Astérisque}, 223:9–57, 1994. §4. \href{https://www.numdam.org/item/AST_1994__223__9_0/}{Original article}.}

\bpbibitem{Kat88}{bp60}{N. M. Katz. Travaux de Laumon. \emph{Astérisque}, 161–162:105–132, 1988. Séminaire Bourbaki, exp. 691.}

\bpbibitem{KW01}{bp61}{R. Kiehl and R. Weissauer. \emph{Weil conjectures, perverse sheaves and \(\ell\)-adic Fourier transform}, volume 42 of \emph{Ergebnisse der Mathematik}. Springer, 2001.}

\bpbibitem{Lau87}{bp62}{G. Laumon. Transformation de Fourier, constantes d'équations fonctionnelles et conjecture de Weil. \emph{Publ. Math. IHÉS}, 65:131–210, 1987.}

\bpbibitem{Mil86a}{bp63}{J. S. Milne. \href{https://www.jmilne.org/math/xnotes/AVs.pdf}{Abelian Varieties}. In \emph{Arithmetic Geometry}. 1986. Corrected online version.}

\bpbibitem{Mil86b}{bp64}{J. S. Milne. \href{https://www.jmilne.org/math/xnotes/JVs.pdf}{Jacobian Varieties}. In \emph{Arithmetic Geometry}. 1986. Corrected online version.}

\bpbibitem{Ran39}{bp65}{R. A. Rankin. Contributions to the theory of Ramanujan's function \(\tau(n)\) and similar arithmetical functions. II. \emph{Proc. Cambridge Philos. Soc.}, 35:357–372, 1939.}

\bpbibitem{SGA4}{bp66}{M. Artin, A. Grothendieck, and J.-L. Verdier, editors. \emph{Théorie des topos et cohomologie étale des schémas}, volumes 269, 270, 305 of \emph{Lecture Notes in Mathematics}. Springer, 1972–1973. SGA 4.}

\bpbibitem{SGA4½}{bp67}{P. Deligne. \emph{Cohomologie étale}, volume 569 of \emph{Lecture Notes in Mathematics}. Springer, 1977. SGA 4½; \emph{Théorèmes de finitude}, §3, and \emph{Rapport sur la formule des traces}, §3, and \emph{Sommes trigonométriques}, §1.}

\bpbibitem{SGA5}{bp68}{A. Grothendieck, editor. \emph{Cohomologie \(\ell\)-adique et fonctions \(L\)}, volume 589 of \emph{Lecture Notes in Mathematics}. Springer, 1977. SGA 5.}

\bpbibitem{SGA7II}{bp69}{P. Deligne and N. Katz, editors. \emph{Groupes de monodromie en géométrie algébrique}, volume 340 of \emph{Lecture Notes in Mathematics}. Springer, 1973. SGA 7 II; Exposé XIII, (2.1.7).}

\bpbibitem{Stacks}{bp70}{The Stacks Project Authors. \href{https://stacks.math.columbia.edu}{The Stacks Project}. Individual results are cited by their stable tags in the text.}

\bpbibitem{Var05}{bp71}{Y. Varshavsky. Lefschetz–Verdier trace formula and a generalization of a theorem of Fujiwara. 2005. Version 2. arXiv:\href{https://arxiv.org/abs/math/0505564}{math/0505564}. \href{https://arxiv.org/abs/math/0505564v2}{Author's preprint}.}

\bpbibitem{Wei48a}{bp72}{A. Weil. \emph{Sur les courbes algébriques et les variétés qui s'en déduisent}. Hermann, Paris, 1948.}

\bpbibitem{Wei48b}{bp73}{A. Weil. \emph{Variétés abéliennes et courbes algébriques}. Hermann, Paris, 1948.}

\bpbibitem{Wei49}{bp74}{A. Weil. Numbers of solutions of equations in finite fields. \emph{Bull. Amer. Math. Soc.}, 55:497–508, 1949.}

\end{document}
